\subsection{Conservation mechanism ablation}\label{app:conservation}
This experiment compares ways of obtaining flow conservation on
Sioux Falls under otherwise identical supervision, asking what the
conservative parameterization of Eq.~\eqref{eq:consparam} contributes
beyond conservation itself:
(i)~\emph{exact} --- the path-share parameterization, pretrained on the
unscaled flow loss (the published five-seed models of
Section~\ref{sec:pairresults});
(ii)~\emph{exact (scaled)} --- the same architecture pretrained on the
per-link scaled loss, the one recipe difference of the earlier
three-way comparison, removed by training the exact parameterization
on the direct models' objective;
(iii)~\emph{penalty} --- a direct MLP on link flows with a squared
node-balance penalty (weight $1$, normalized by the nodal-throughput
RMS over the training states, floored at the network average);
(iv)~\emph{projection} --- the same direct MLP whose output is
projected each step onto $\{f:Bf=Mq+b\}$ by the differentiable
least-norm map $f\mapsto f-B^{\!\top}(BB^{\!\top})^{\dagger}(Bf-Mq-b)$
(Moore--Penrose inverse; $BB^{\!\top}$ is singular for a connected
network). The direct-MLP width is solved from the parameter-count
equation so that all four models have $714{,}979$ parameters against
the exact model's $715{,}040$. Pretraining shares the data, optimizer
(Adam $3\times10^{-3}$, batch 256, at most 12{,}000 steps) and the
min-validation checkpoint rule (validation metric: unscaled relative
L2 flow error) across all architectures; the loss differs only as
stated. Continuation follows the protocol of
Appendix~\ref{app:pairs-v102} unchanged (1{,}200 updates, shared
labels and minibatch sequences, the uniform selection rule including
the pre-trained initial model in the candidate pool), with flow-only
and response supervision for each mechanism. The exact rows reproduce
the published five-seed medians exactly ($21.18\%$ and $14.19\%$
column error), and every projection output satisfies $Bf=Mq+b$ to
$2\times10^{-11}$; the path-share construction holds node balance to
the same scale under both pretraining losses.

\begin{table}[pos=htbp]\centering\footnotesize
\caption{Conservation mechanism ablation on the 96 Sioux Falls test
paths (medians of path-level values; active queries and seeds averaged
within path). Mechanisms: path-share parameterization under the two
pretraining losses, and the two direct mechanisms under the scaled
loss. Residuals are per-query maxima over the path states, maximized
over queries: the node-balance residual $\max_q\|Bf(q)-Mq-b\|_\infty$
(veh/h); the derivative residual
$\max_{q,j}\|\partial_j[Bf](q)-\mu_j\|_\infty/\|\mu_j\|_2$. Column:
attribution column error; incremental: endpoint flow-change error;
$E_{\mathrm{total}}$: normalized end-to-end summary error.}\label{tab:conservation}
\setlength{\tabcolsep}{3.6pt}
\begin{tabular*}{\tblwidth}{@{\extracolsep{\fill}}llrrrrrr}
\toprule
Mechanism & Method & Flow (\%) & Incremental (\%) & Node resid. & Deriv.\ resid. & Column (\%) & $E_{\mathrm{total}}$\\
\midrule
Exact & flow & 0.032 & 5.03 & $3\times10^{-11}$ & $9\times10^{-16}$ & 21.18 & 11.35\\
Exact & response & 0.029 & 3.52 & $3\times10^{-11}$ & $8\times10^{-16}$ & 14.19 & 8.68\\
\addlinespace
Exact (scaled) & flow & 0.033 & 5.53 & $3\times10^{-11}$ & $8\times10^{-16}$ & 21.60 & 10.33\\
Exact (scaled) & response & 0.032 & 3.81 & $3\times10^{-11}$ & $8\times10^{-16}$ & 14.35 & 8.19\\
\addlinespace
Penalty & flow & 0.040 & 22.76 & 42.1 & 1.16 & 183.7 & 148.8\\
Penalty & response & 0.039 & 22.24 & 42.1 & 1.17 & 177.9 & 144.9\\
\addlinespace
Projection & flow & 0.026 & 14.54 & $2\times10^{-11}$ & $2\times10^{-15}$ & 104.1 & 46.4\\
Projection & response & 0.026 & 14.42 & $2\times10^{-11}$ & $2\times10^{-15}$ & 99.4 & 45.6\\
\bottomrule\end{tabular*}\end{table}

Table~\ref{tab:conservation} supports three readings. First, the
pretraining loss is not what separates the exact parameterization from
the direct mechanisms: under the same scaled loss the exact
architecture attains $21.60\%$ and $14.35\%$ column error against the
projection's $104.1\%$ and $99.4\%$, and training the exact
architecture on the scaled loss leaves its column errors within half a
percentage point of the published unscaled values ($21.18\%$,
$14.19\%$) while slightly lowering its end-to-end summary error. The
exact-versus-direct gap is therefore a property of the
parameterization, not of its training recipe. Second, conservation
itself still matters: the penalty model matches the flow accuracy of
the exact model ($0.040\%$) with a small output residual ($42$ veh/h
against nodal throughputs in the thousands), yet its integrated
attribution residual --- the quantity $R_j$ of
Eq.~\eqref{eq:residualidentity} --- reaches $0.66$ in units of the
demand-conservation scale $|\Delta q_j|\|\mu_j\|_2$, its derivative
balance residual is $1.16$ in the same normalization, and the
downstream outcome columns inherit the damage; the chain of
Eq.~\eqref{eq:residualidentity} is exactly what fails, because the
penalty controls $r(q)$ but not $\partial_jr$. The projection enforces
output conservation exactly and transfers it to derivatives and to the
integrated residual, cutting the column error to $104.1\%$ while
remaining far above the exact parameterization, because projecting an
unconstrained response field onto node balance does not reproduce the
path-share structure that the assignment imposes. Third, response
supervision improves every mechanism (penalty $183.7\to177.9$,
projection $104.1\to99.4$, exact $21.18\to14.19$ and
$21.60\to14.35$), with by far the largest effect inside the exact
parameterization: the inductive benefit of
conservation-by-construction and the statistical benefit of derivative
supervision are complementary rather than substitutes.
